\chapter{Regression}
\label{ch:regression}

\epigraph{``All of machine learning is either regression or classification. Everything else is a trick.''}{}

\learninggoal{
	\item Derive and interpret linear regression via ordinary least squares.
	\item Understand the assumptions underlying linear regression and how violations affect inference.
	\item Extend linear regression with basis functions for non-linear relationships.
	\item Diagnose model fit using residual analysis and influence measures.
}

\section{Linear Regression}

The workhorse of supervised learning: model the target as a linear combination of features.

\[
	y = \mathbf{w}^\top \mathbf{x} + \epsilon, \quad \epsilon \sim \N(0, \sigma^2)
\]

Given data $\{(\mathbf{x}_i, y_i)\}_{i=1}^n$, collect inputs into matrix $\mathbf{X} \in \R^{n \times d}$ and outputs into vector $\mathbf{y} \in \R^n$:

\[
	\mathbf{X} = \begin{pmatrix}
		\mathbf{x}_1^\top \\
		\vdots            \\
		\mathbf{x}_n^\top
	\end{pmatrix}, \quad
	\mathbf{y} = \begin{pmatrix} y_1 \\ \vdots \\ y_n \end{pmatrix}
\]

\subsection{Ordinary Least Squares}

The OLS estimator minimizes the sum of squared residuals:

\[
	\hat{\mathbf{w}}_{\text{OLS}} = \arg\min_{\mathbf{w}} \|\mathbf{X}\mathbf{w} - \mathbf{y}\|_2^2
\]

Setting the gradient to zero gives the \emph{normal equations}:

\[
	\mathbf{X}^\top \mathbf{X} \mathbf{w} = \mathbf{X}^\top \mathbf{y}
\]
\[
	\hat{\mathbf{w}}_{\text{OLS}} = (\mathbf{X}^\top \mathbf{X})^{-1} \mathbf{X}^\top \mathbf{y}
\]

\begin{definition}[Fitted Values and Residuals]
	\textbf{Fitted values:} $\hat{\mathbf{y}} = \mathbf{X} \hat{\mathbf{w}} = \mathbf{H} \mathbf{y}$, where $\mathbf{H} = \mathbf{X}(\mathbf{X}^\top \mathbf{X})^{-1} \mathbf{X}^\top$ is the \textbf{hat matrix}.\\
	\textbf{Residuals:} $\mathbf{e} = \mathbf{y} - \hat{\mathbf{y}}$.
\end{definition}

\subsection{Geometric Interpretation}

OLS projects $\mathbf{y}$ onto the column space of $\mathbf{X}$. The hat matrix $\mathbf{H}$ is the projection matrix. The residuals $\mathbf{e}$ are orthogonal to $\text{col}(\mathbf{X})$:

\[
	\mathbf{X}^\top \mathbf{e} = \mathbf{0}
\]

\section{Assumptions and Diagnostics}

Linear regression makes four key assumptions:

\begin{longtable}{p{2.5cm}p{5cm}p{5cm}}
	\toprule
	\textbf{Assumption} & \textbf{Meaning}                                       & \textbf{Diagnostic}        \\
	\midrule
	Linearity           & $\E[y \mid \mathbf{x}]$ is linear in $\mathbf{x}$      & Residuals vs.\ fitted plot \\
	Independence        & Observations are independent                           & Durbin--Watson test        \\
	Homoscedasticity    & $\Var(\epsilon \mid \mathbf{x}) = \sigma^2$ (constant) & Scale--location plot       \\
	Normality           & $\epsilon \sim \N(0, \sigma^2)$                        & Q--Q plot of residuals     \\
	\bottomrule
\end{longtable}

\subsection{Influential Points}

A point is \emph{influential} if removing it substantially changes $\hat{\mathbf{w}}$. Measure with Cook's distance:

\[
	D_i = \frac{e_i^2}{d \cdot \text{MSE}} \cdot \frac{h_{ii}}{(1 - h_{ii})^2}
\]

where $h_{ii}$ is the $i$-th diagonal of $\mathbf{H}$ (the \emph{leverage}). A common threshold: $D_i > 4/n$ warrants investigation.

\section{Basis Function Expansion}

Linear regression is limited to linear relationships---but only the \emph{parameters} must be linear. The features can be non-linear transformations:

\[
	y = w_0 + w_1 \phi_1(x) + w_2 \phi_2(x) + \cdots + w_k \phi_k(x) + \epsilon
\]

\subsection{Polynomial Basis}

\[
	\phi_j(x) = x^j
\]

\begin{lstlisting}[language=Python, caption=Polynomial regression]
from sklearn.preprocessing import PolynomialFeatures
from sklearn.linear_model import LinearRegression

poly = PolynomialFeatures(degree=3)
X_poly = poly.fit_transform(X)
model = LinearRegression().fit(X_poly, y)
\end{lstlisting}

Choosing the polynomial degree is a bias--variance tradeoff. Low degrees underfit; high degrees overfit.

\subsection{Spline Basis}

Splines are piecewise polynomials joined at \emph{knots}. Cubic splines are the most common: each piece is a cubic polynomial, with continuity constraints at knots.

\[
	y = \beta_0 + \beta_1 x + \beta_2 x^2 + \beta_3 x^3 + \sum_{k=1}^K \theta_k (x - \xi_k)_+^3 + \epsilon
\]

where $(z)_+ = \max(0, z)$ and $\xi_k$ are the knot locations. Splines are more stable than high-degree polynomials.

\subsection{Regression Splines vs. Smoothing Splines}

\begin{itemize}
	\item \textbf{Regression splines:} choose knot locations and basis functions, then fit by OLS.
	\item \textbf{Smoothing splines:} place a knot at every data point and add a penalty on the second derivative to prevent overfitting.
\end{itemize}

\section{Model Selection in Regression}

\subsection{Adjusted $R^2$}

\[
	R^2_{\text{adj}} = 1 - \frac{\text{RSS} / (n - d - 1)}{\text{TSS} / (n - 1)}
\]

Penalizes adding useless features, unlike plain $R^2$.

\subsection{Information Criteria}

\begin{itemize}
	\item \textbf{AIC:} $-2 \log L + 2d$
	\item \textbf{BIC:} $-2 \log L + d \log n$
\end{itemize}

Both balance fit (log-likelihood) against model complexity. BIC penalizes complexity more heavily (useful for selecting the true model when it exists). AIC is better for prediction.

\section{Multiple Regression and Interaction Effects}

When features interact, their joint effect is not the sum of individual effects:

\[
	y = w_0 + w_1 x_1 + w_2 x_2 + w_{12} x_1 x_2 + \epsilon
\]

The interaction term $x_1 x_2$ allows the effect of $x_1$ to depend on $x_2$:

\[
	\frac{\partial y}{\partial x_1} = w_1 + w_{12} x_2
\]

Including all pairwise interactions adds $d(d-1)/2$ terms. Feature engineering often involves discovering important interactions.

\section{Exercises}

\begin{exercise}
	Derive the normal equations by setting the gradient of $\|\mathbf{X}\mathbf{w} - \mathbf{y}\|_2^2$ to zero. Show each step.
\end{exercise}

\begin{exercise}
	A dataset has 3 features and 100 observations. You fit a linear regression and obtain $R^2 = 0.85$ and $R^2_{\text{adj}} = 0.84$. You add 10 random noise features and refit. What do you expect to happen to $R^2$ and $R^2_{\text{adj}}$?
\end{exercise}

\begin{exercise}
	Generate data from $y = \sin(x) + \epsilon$ with $x \in [0, 2\pi]$. Fit polynomials of degree 1, 3, 5, 15. Plot the fitted curves and report the test MSE on a held-out set. At what degree does overfitting begin?
\end{exercise}

\begin{exercise}
	Explain why using a high-degree polynomial can produce wild extrapolations beyond the data range, while cubic splines are better behaved. What property of polynomials causes this?
\end{exercise}
